% Presentation of original arrow notation. See flow-manifest.json.
\begin{center}
\begin{tikzpicture}[book flow node/.style={draw=BookHeading,line width=.65pt,fill=BookPanel,rounded corners=1.5pt,align=center,inner sep=5pt,minimum height=10mm,font=\fontsize{9}{11}\selectfont,execute at begin node={\hyphenpenalty=10000\exhyphenpenalty=10000}},book flow edge/.style={-{Latex[length=2mm]},draw=BookHeading,line width=.7pt}]
\node[book flow node,text width=49mm] (n0) at (0,0) {Theme: id, label};
\node[book flow node,text width=49mm] (n1) at (-3,-2.5) {Mode: dark /\allowbreak{} token to color};
\node[book flow node,text width=49mm] (n2) at (3,-2.5) {Mode: light /\allowbreak{} token to color};
\node[book flow node,text width=49mm] (n3) at (0,-5) {Style: fg token\allowbreak{}Ref(.\allowbreak{}.\allowbreak{}.\allowbreak{}), bg token\allowbreak{}Ref(.\allowbreak{}.\allowbreak{}.\allowbreak{}), layout rules};
\node[book flow node,text width=49mm] (n4) at (0,-7.5) {Surface cells: concrete colors};
\draw[book flow edge] (n0) -- node[fill=white,inner sep=2pt,align=center,font=\scriptsize] {owns mode maps} (n1);
\draw[book flow edge] (n0) -- node[fill=white,inner sep=2pt,align=center,font=\scriptsize] {owns mode maps} (n2);
\draw[book flow edge] (n1) -- node[fill=white,inner sep=2pt,align=center,font=\scriptsize] {resolves refs} (n3);
\draw[book flow edge] (n2) -- node[fill=white,inner sep=2pt,align=center,font=\scriptsize] {resolves refs} (n3);
\draw[book flow edge] (n3) -- node[fill=white,inner sep=2pt,align=center,font=\scriptsize] {render context} (n4);
\end{tikzpicture}
\end{center}
